\documentclass[12pt,a4paper]{article}

% =====================================================
% PAKETLER
% =====================================================
\usepackage[utf8]{inputenc}
\usepackage[T1]{fontenc}
\usepackage[turkish]{babel}
\usepackage{geometry}
\usepackage{graphicx}
\usepackage{booktabs}
\usepackage{array}
\usepackage{amsmath}
\usepackage{hyperref}
\usepackage{float}
\usepackage{fancyhdr}
\usepackage{titlesec}
\usepackage{xcolor}
\usepackage{listings}
\usepackage{caption}
\usepackage{subcaption}
\usepackage{multirow}
\usepackage{setspace}

% =====================================================
% SAYFA DÜZENİ
% =====================================================
\geometry{
    top=2.5cm, bottom=2.5cm,
    left=3cm,  right=2.5cm
}
\onehalfspacing

% =====================================================
% BAŞLIK STİLİ
% =====================================================
\titleformat{\section}{\large\bfseries}{}{0em}{\thesection.\quad}
\titleformat{\subsection}{\normalsize\bfseries}{}{0em}{\thesubsection.\quad}

% =====================================================
% ÜSTBILGI / ALTBILGI
% =====================================================
\pagestyle{fancy}
\fancyhf{}
\rhead{\small NLP Projesi -- Chunking}
\lhead{\small Bursa Teknik Üniversitesi}
\cfoot{\thepage}
\renewcommand{\headrulewidth}{0.4pt}

% =====================================================
% KOD STİLİ
% =====================================================
\lstset{
    language=Python,
    basicstyle=\ttfamily\small,
    keywordstyle=\color{blue}\bfseries,
    commentstyle=\color{gray},
    stringstyle=\color{orange!80!black},
    frame=single,
    breaklines=true,
    numbers=left,
    numberstyle=\tiny\color{gray},
    showstringspaces=false,
    captionpos=b,
}

% =====================================================
% BELGE BAŞLADI
% =====================================================
\begin{document}

% -------------------------------------------
% KAPAK SAYFASI
% -------------------------------------------
\begin{titlepage}
    \centering
    \vspace*{1cm}

    {\Large \textbf{T.C.}\\[0.3em]
     \textbf{BURSA TEKNİK ÜNİVERSİTESİ}\\[0.3em]
     \textbf{Mühendislik ve Doğa Bilimleri Fakültesi}\\[0.3em]
     \textbf{Bilgisayar Mühendisliği Bölümü}}

    \vspace{1.5cm}
    \vspace{2cm}

    \rule{\linewidth}{1pt}\\[0.5em]
    {\LARGE \textbf{Doğal Dil İşleme Projesi}}\\[0.5em]
    {\Large \textbf{İsim ve Diğer Öbeklerin Saptanması (Chunking)}}\\[0.3em]
    \rule{\linewidth}{1pt}

    \vspace{1.5cm}

    \begin{tabular}{ll}
        \textbf{Ders:}    & Doğal Dil İşleme \\[0.5em]
        \textbf{Dönem:}   & 2025 -- 2026 Bahar Dönemi \\[0.5em]
        \textbf{Öğrenci:} & [AD SOYAD] \\[0.5em]
        \textbf{Numara:}  & [ÖĞRENCİ NUMARASI] \\[0.5em]
        \textbf{Grup Üyeleri:} & [VARSA DİĞER ÜYELERİN ADI SOYADI] \\[0.5em]
        \textbf{Danışman:}& [HOCA ADI] \\
    \end{tabular}

    \vfill
    {\large \today}
\end{titlepage}


% =====================================================
% 1. GİRİŞ
% =====================================================
\section{Giriş}

Bu proje, Bursa Teknik Üniversitesi Bilgisayar Mühendisliği Bölümü
2025--2026 Bahar Dönemi Doğal Dil İşleme dersi kapsamında hazırlanmıştır.

Projenin konusu \textbf{İsim ve Diğer Öbeklerin Saptanması (Chunking)} olarak belirlenmiştir.
Chunking (ya da kısmi ayrıştırma), bir cümle içindeki kelimelerin hangi sözdizimsel gruba
(öbeğe) ait olduğunu belirleyen bir Doğal Dil İşleme görevidir. Bu görevde:

\begin{itemize}
    \item \textbf{NP} (Noun Phrase / İsim Öbeği): İsimlerin oluşturduğu gruplar
    \item \textbf{VP} (Verb Phrase / Eylem Öbeği): Fiillerin oluşturduğu gruplar
    \item \textbf{ADVP} (Adverb Phrase / Zarf Öbeği): Zarfların oluşturduğu gruplar
    \item \textbf{PP} (Prepositional Phrase / Edat Öbeği): Edatların oluşturduğu gruplar
    \item \textbf{O}: Herhangi bir öbeğe dahil olmayan kelimeler
\end{itemize}

etiketleri BIO (Beginning--Inside--Outside) formatında her kelimeye atanmaktadır.

\subsection{BIO Etiketleme Formatı}

BIO formatında her kelime üç durumdan birinde etiketlenir:
\begin{itemize}
    \item \textbf{B-X}: Bir X öbeğinin \textit{başlangıç} kelimesi
    \item \textbf{I-X}: Bir X öbeğinin \textit{devam} kelimesi
    \item \textbf{O}: Öbek dışı kelime
\end{itemize}

Örnek cümle ve CoNLL formatındaki etiketleme:

\begin{table}[H]
\centering
\caption{CoNLL formatında örnek etiketleme}
\begin{tabular}{cllll}
\toprule
\textbf{ID} & \textbf{FORM} & \textbf{CHUNK-OUTER} & \textbf{CHUNK-INNER} & \textbf{CLAUSE} \\
\midrule
1  & Dün         & B-ADVP & \_ & O \\
2  & akşam       & I-ADVP & \_ & O \\
3  & toplantıdan & B-NP   & B-RELCL & B-RELCL \\
4  & erken       & I-NP   & I-RELCL & I-RELCL \\
5  & çıkan       & I-NP   & I-RELCL & I-RELCL \\
6  & öğrencinin  & I-NP   & \_ & O \\
7  & ,           & O      & \_ & O \\
8  & hocasının   & B-NP   & B-RELCL & B-RELCL \\
9  & önerdiği    & I-NP   & I-RELCL & I-RELCL \\
10 & makaleyi    & I-NP   & \_ & O \\
11 & kütüphanede & B-NP   & \_ & B-COMPCL \\
12 & dikkatlice  & B-ADVP & \_ & I-COMPCL \\
13 & okuduğunu   & B-VP   & \_ & I-COMPCL \\
14 & fark        & B-VP   & \_ & O \\
15 & ettim       & I-VP   & \_ & O \\
16 & .           & O      & \_ & O \\
\bottomrule
\end{tabular}
\end{table}

% =====================================================
% 2. VERİ SETİ
% =====================================================
\section{Veri Seti}

\subsection{Universal Dependencies Türkçe Treebank}

Bu projede veri seti olarak \textbf{Universal Dependencies (UD) Turkish-BOUN Treebank}
kullanılmıştır. Bu veri seti, Boğaziçi Üniversitesi tarafından hazırlanmış olup
CoNLL-U formatında ayrıntılı Türkçe sözdizimsel etiketlemeler içermektedir
\cite{ud_turkish}.

Veri setinin temel özellikleri:

\begin{table}[H]
\centering
\caption{Veri seti istatistikleri}
\begin{tabular}{lrr}
\toprule
\textbf{Bölüm} & \textbf{Cümle Sayısı} & \textbf{Token Sayısı} \\
\midrule
Eğitim (train) & 8.782 & 113.002 \\
Test           &   979 &  12.210 \\
\textbf{Toplam} & \textbf{9.761} & \textbf{125.212} \\
\bottomrule
\end{tabular}
\end{table}

\subsection{Veri Hazırlama}

UD Treebank, CoNLL-U formatında UPOS (Universal Part-of-Speech) etiketleri içermektedir.
Bu etiketler aşağıdaki kurala göre chunk etiketlerine dönüştürülmüştür:

\begin{table}[H]
\centering
\caption{POS etiketi $\rightarrow$ Chunk etiketi dönüşüm kuralları}
\begin{tabular}{ll}
\toprule
\textbf{UPOS Etiketleri} & \textbf{Chunk Türü} \\
\midrule
NOUN, PROPN, PRON, DET, ADJ, NUM & NP (İsim Öbeği) \\
VERB, AUX                         & VP (Eylem Öbeği) \\
ADV                               & ADVP (Zarf Öbeği) \\
ADP                               & PP (Edat Öbeği) \\
PUNCT, CONJ, CCONJ, SCONJ, diğer & O \\
\bottomrule
\end{tabular}
\end{table}

\subsection{Etiket Dağılımı}

Eğitim verisindeki etiket dağılımı aşağıdaki tabloda verilmiştir:

\begin{table}[H]
\centering
\caption{Eğitim verisi etiket dağılımı}
\begin{tabular}{lrr}
\toprule
\textbf{Etiket} & \textbf{Token Sayısı} & \textbf{Oran (\%)} \\
\midrule
B-NP   & 27.690 & 24.50 \\
I-NP   & 30.885 & 27.33 \\
B-VP   & 18.216 & 16.12 \\
I-VP   &  4.588 &  4.06 \\
B-ADVP &  4.405 &  3.90 \\
I-ADVP &    527 &  0.47 \\
B-PP   &  2.222 &  1.97 \\
I-PP   &     11 &  0.01 \\
O      & 24.458 & 21.64 \\
\midrule
\textbf{Toplam} & \textbf{113.002} & \textbf{100.00} \\
\bottomrule
\end{tabular}
\end{table}

% =====================================================
% 3. YÖNTEM
% =====================================================
\section{Yöntem}

\subsection{Öznitelik Mühendisliği}

Her kelime için aşağıdaki öznitelikler çıkarılmıştır:

\begin{itemize}
    \item \textbf{Kelime öznitelikleri:} Kelimenin kendisi, küçük harf formu, büyük harf mi, baş harfi büyük mü, rakam mı, uzunluğu
    \item \textbf{Suffix (son ekler):} Son 1, 2, 3 ve 4 karakter -- Türkçe eklemeli bir dil olduğundan morfolojik eklere dayalı bu öznitelikler çok kritiktir
    \item \textbf{Prefix (ön ekler):} İlk 2 ve 3 karakter
    \item \textbf{Bağlam penceresi:} Önceki 2 ve sonraki 2 kelime ile bunların özellikleri
    \item \textbf{Cümle sınırları:} BOS (cümle başı), EOS (cümle sonu) bayrakları
\end{itemize}

\subsection{Model 1: Conditional Random Fields (CRF)}

\textbf{Conditional Random Fields (Koşullu Rastsal Alanlar)}, dizi etiketleme
problemleri için tasarlanmış istatistiksel bir modeldir. Naive Bayes ve Logistic
Regression'dan farklı olarak her kelimeyi bağımsız sınıflandırmak yerine,
komşu etiketler arasındaki bağımlılıkları da modelleyebilir. Bu özelliği
chunking gibi ardışık etiketleme görevleri için CRF'yi ideal kılmaktadır.

Model parametreleri:
\begin{itemize}
    \item Algoritma: L-BFGS (Limited-memory Broyden--Fletcher--Goldfarb--Shanno)
    \item L1 düzenlileştirme katsayısı: $c_1 = 0.05$
    \item L2 düzenlileştirme katsayısı: $c_2 = 0.05$
    \item Maksimum iterasyon: 150
\end{itemize}

\subsection{Model 2: Naive Bayes}

\textbf{Multinomial Naive Bayes}, Bayes teoremini temel alan istatistiksel bir sınıflandırıcıdır.
Özellikler arasındaki bağımsızlık varsayımı yapsa da hızlı ve basit yapısıyla
temel (baseline) model olarak kullanılmaktadır. CRF ile karşılaştırma amacıyla
ikinci model olarak seçilmiştir.

Model parametreleri:
\begin{itemize}
    \item Düzleştirme katsayısı: $\alpha = 0.1$
    \item Öznitelik ölçekleme: MinMaxScaler (Naive Bayes negatif değer kabul etmez)
\end{itemize}

% =====================================================
% 4. SONUÇLAR
% =====================================================
\section{Sonuçlar}

\subsection{Model Karşılaştırması}

\begin{table}[H]
\centering
\caption{Model başarı karşılaştırması}
\begin{tabular}{lccc}
\toprule
\textbf{Model} & \textbf{Accuracy} & \textbf{Macro F1} & \textbf{Weighted F1} \\
\midrule
CRF         & \textbf{0.9283} & \textbf{0.7755} & \textbf{0.9278} \\
Naive Bayes & 0.5129          & 0.2019          & 0.4357          \\
\bottomrule
\end{tabular}
\end{table}

CRF modeli, chunking gibi sıralı bağımlılıkların kritik önem taşıdığı bir görevde
\textbf{\%92.83 doğruluk oranı} elde etmiştir. Naive Bayes ise her kelimeyi
bağımsız değerlendirdiğinden \%51.29 ile sınırlı kalmıştır. Bu fark,
komşu etiket bilgisinin chunking başarısındaki kritik rolünü açıkça göstermektedir.

\subsection{CRF Sınıf Bazlı Metrikler}

\begin{table}[H]
\centering
\caption{CRF modeli -- sınıf bazlı sonuçlar}
\begin{tabular}{lrrrr}
\toprule
\textbf{Etiket} & \textbf{Precision} & \textbf{Recall} & \textbf{F1-Score} & \textbf{Destek} \\
\midrule
B-ADVP & 0.8212 & 0.8097 & 0.8154 &   431 \\
B-NP   & 0.9245 & 0.9316 & 0.9280 & 2.996 \\
B-PP   & 0.9528 & 0.8441 & 0.8952 &   263 \\
B-VP   & 0.9100 & 0.8784 & 0.8939 & 2.048 \\
I-ADVP & 0.7692 & 0.6122 & 0.6818 &    49 \\
I-NP   & 0.9186 & 0.9497 & 0.9338 & 3.456 \\
I-PP   & 0.0000 & 0.0000 & 0.0000 &     1 \\
I-VP   & 0.8699 & 0.8210 & 0.8447 &   391 \\
O      & 0.9864 & 0.9864 & 0.9864 & 2.575 \\
\midrule
\textbf{Accuracy}     & \multicolumn{3}{c}{0.9283} & 12.210 \\
\textbf{Macro avg}    & 0.7947 & 0.7592 & 0.7755 & 12.210 \\
\textbf{Weighted avg} & 0.9279 & 0.9283 & 0.9278 & 12.210 \\
\bottomrule
\end{tabular}
\end{table}

\textit{Not: I-PP sınıfı test setinde yalnızca 1 örnekle temsil edildiğinden sıfır F1 skoru elde edilmiştir. Bu istatistiksel olarak anlamlı değildir.}

\subsection{Grafiksel Sonuçlar}

\subsubsection{CRF Karmaşıklık Matrisi}

\begin{figure}[H]
    \centering
    \includegraphics[width=0.90\textwidth]{confusion_matrix_crf.png}
    \caption{CRF modeli karmaşıklık matrisi (Confusion Matrix). Köşegen üzerindeki yüksek değerler modelin başarılı sınıflandırma yaptığını göstermektedir.}
    \label{fig:cm_crf}
\end{figure}

\subsubsection{CRF Sınıf Bazlı Metrik Grafiği}

\begin{figure}[H]
    \centering
    \includegraphics[width=\textwidth]{metrics_crf.png}
    \caption{CRF modeli -- her sınıf için Precision, Recall ve F1-Score değerleri.}
    \label{fig:metrics_crf}
\end{figure}

\subsubsection{Naive Bayes Karmaşıklık Matrisi}

\begin{figure}[H]
    \centering
    \includegraphics[width=0.90\textwidth]{confusion_matrix_nb.png}
    \caption{Naive Bayes modeli karmaşıklık matrisi. Köşegen dışı değerlerin yüksekliği modelin birçok sınıfı birbirine karıştırdığını göstermektedir.}
    \label{fig:cm_nb}
\end{figure}

\subsubsection{Model Karşılaştırma Grafiği}

\begin{figure}[H]
    \centering
    \includegraphics[width=0.80\textwidth]{model_comparison.png}
    \caption{CRF ve Naive Bayes modellerinin Accuracy, Macro F1 ve Weighted F1 skorları karşılaştırması.}
    \label{fig:comparison}
\end{figure}

\subsection{Tartışma}

CRF modeli tüm metriklerde Naive Bayes'i önemli ölçüde geride bırakmıştır.
Bu sonucun temel nedenleri şöyle sıralanabilir:

\begin{enumerate}
    \item \textbf{Dizi bağımlılığı:} CRF, bir kelimenin etiketini belirlerken komşu etiketleri
    de göz önünde bulundurur. Chunking'de bir NP öbeği içindeki kelimelerin birbirini
    takip etmesi gibi dizi kalıpları bu sayede yakalanabilmektedir.

    \item \textbf{Türkçe'ye uygunluk:} Türkçe eklemeli bir dildir; kelimenin sonu çoğu zaman
    sözdizimsel rolünü belirler. CRF'nin suffix özniteliklerini etkili biçimde kullanması
    Türkçe'de özellikle faydalı olmuştur.

    \item \textbf{Naive Bayes'in sınırlılıkları:} Naive Bayes bağımsızlık varsayımı nedeniyle
    dizi içi bağımlılıkları modelleyemez. Bu nedenle özellikle B-VP, I-VP, I-ADVP gibi
    ardışıklık gerektiren sınıflarda başarısız olmuştur.
\end{enumerate}

% =====================================================
% 5. SONUÇ
% =====================================================
\section{Sonuç}

Bu projede Türkçe metinler için chunking görevi gerçekleştirilmiştir. Universal Dependencies
Turkish-BOUN Treebank kullanılarak 9.761 cümlelik bir veri seti hazırlanmıştır. İki farklı
makine öğrenmesi modeli eğitilmiş ve karşılaştırılmıştır:

\begin{itemize}
    \item \textbf{CRF modeli \%92.83 doğruluk oranı} elde ederek chunking için güçlü bir
    performans sergilemiştir. Özellikle B-NP (\%93 F1), I-NP (\%93 F1) ve O (\%99 F1)
    sınıflarında çok başarılı sonuçlar alınmıştır.

    \item \textbf{Naive Bayes modeli \%51.29 doğruluk oranı} ile kalmıştır. Dizi
    bağımlılıklarını yakalayamaması bu düşük performansın temel nedenidir.
\end{itemize}

Sonuçlar, dizi etiketleme problemlerinde CRF gibi bağlam duyarlı modellerin
klasik istatistiksel modellere olan üstünlüğünü açıkça ortaya koymaktadır.

% =====================================================
% KAYNAKÇA
% =====================================================
\begin{thebibliography}{9}

\bibitem{ud_turkish}
Türk, U., Atmaca, F., Özateş, Ş. B., Köse, A., Küçükçiloğlu, S., Eryiğit, G.,
\& Özgür, A. (2022).
\textit{Resources for Turkish Dependency Parsing: Introducing the BOUN Treebank
and the BoAT Annotation Tool}.
Language Resources and Evaluation, 56, 579--625.

\bibitem{lafferty2001crf}
Lafferty, J., McCallum, A., \& Pereira, F. (2001).
\textit{Conditional Random Fields: Probabilistic Models for Segmenting and
Labeling Sequence Data}.
In Proceedings of ICML, 282--289.

\bibitem{jurafsky2023slp}
Jurafsky, D., \& Martin, J. H. (2023).
\textit{Speech and Language Processing} (3rd ed. draft).
\url{https://web.stanford.edu/~jurafsky/slp3/}

\end{thebibliography}

\end{document}
